\section{Background Literature}\label{backg}
% talk about latex3 stuff if there's time and space
\subsection{Choice of programming language}
There were many considerations when deciding what programming language to write \lintexspaced in. Languages chosen were based on those that I had prior experience in, to ensure minimal hinderance from the language itself to the development of the project. An important factor is ensuring \lintex's maintainability in the future. Therefore, only the languages which have actively maintained \acrfull{peg} libraries will be considered, as can be seen in \tableref{table:language_peg_frameworks}. The most recent update times provided in \tableref{table:language_peg_frameworks} are in reference to 24\textsuperscript{th} November 2025.

\begin{table}[h]
	\caption{Languages and related \acrshort{peg} frameworks}
	\label{table:language_peg_frameworks}
	\centering
	\resizebox{\textwidth}{!}{\begin{tabular}{ccc}
		Language & \acrshort{peg} library link & \acrshort{peg} library is in development \\\hline
		Python & https://pythonhosted.org/pyPEG2/ & No, last updated in October 2015 \\
		Python & https://github.com/erikrose/parsimonious & Yes, last updated in November 2025 \\
		JavaScript & https://github.com/pegjs/pegjs & No, last updated in November 2019 \\
		JavaScript & htps://peggyjs.org/ & Yes, last updated in September 2025 \\
		Go & https://github.com/pointlander/peg & Yes, last updated in November 2025 \\
		Go & https://github.com/mna/pigeon & Yes, last updated in November 2025 \\
		C & N/A & N/A \\
		Java & https://github.com/sirthias/parboiled & Yes, last updated in November 2025 \\
		C\# & https://code.google.com/archive/p/peg-sharp/ & No, last updated in May 2011 \\
		C\# & https://github.com/otac0n/Pegasus & No, last updated in December 2023 \\ 
		C++ & https://github.com/yhirose/cpp-peglib & Yes, last updated in July 2025 \\
		PHP & https://github.com/hafriedlander/php-peg & No, last updated in May 2014 \\
		TypeScript & https://github.com/metadevpro/ts-pegjs & No, last updated in November 2024 \\
		TypeScript & https://github.com/EoinDavey/tsPEG & Yes, last updated in November 2025 \\
		Dart & https://github.com/darmie/dart-peg & No, last updated in July 2015 \\
	\end{tabular}}
\end{table}

Based on the current criteria, we can rule out C, C\#, PHP, and Dart.

As many languages remain, there is a need for a stricter criteria: suitability to be utilised as a web application as well as compatibility with the University of Strathclyde's development web server (devweb) system.

\begin{table}[h]
	\caption{Languages and related website development support}
	\label{table:language_web_dev}
    \centering
    \begin{tabular}{ccc}
    Language & Web Development Support & DevWeb Support \\\hline
    Python & Yes & Yes \\
    JavaScript & Yes & Yes \\
    Go & Yes & No \\
    C & No, except for rare exceptions & No \\
    Java & Yes, but not generally used & Yes \\
    C\# & Yes, but not generally used & No \\ 
    C++ & Yes, but not generally used & No \\
    PHP & Yes & Yes \\
    TypeScript & Yes & Yes \\
    Dart & Yes & Yes \\
    \end{tabular}
\end{table}

It should be noted that while Go, C, C\#, and C++ are not directly compatible with devweb, as referenced in \tableref{table:language_web_dev}, they can be supported through use of \acrfull{cgi} scripts, and consequently should not be fully disregarded for this alone. However, based on their support for web development, C, Java, C\#, and C++ will have a lower priority for being a chosen language.

This leaves us with a focus placed on JavaScript, Go, and TypeScript. However, JavaScript will be removed from consideration as a strictly typed language is preferred for our context. Therefore, leaving Go and TypeScript as the main considerations for this project. Finally, as I have more experience and a stronger preference towards Go as opposed to TypeScript, it was the chosen programming language for the project.

\subsection{\nmu Choosing \nmu a \acrshort{peg} library under the Go programming language}\label{backg:peg_library_choice}
Go has two libraries that are used in relation to \acrshort{peg} systems: Pigeon by Martin Angers (mna), and \acrshort{peg} by Andrew Snodgrass (pointlander). Direct comparisons with core components of both have been made in \tableref{table:peg_pigeon_comparison}.

\begin{table}[h]
	\caption{Comparison between PEG and Pigeon}
	\label{table:peg_pigeon_comparison}
	\centering
	\resizebox{\textwidth}{!}{\begin{tabular}{cp{70mm}p{70mm}}
		& Pigeon & Peg \\\hline
		Installation & Requires running a \lstinline[language=Go]!go install! command and setting up environment variables & Requires supporting bash scripting (through Window's Subsystem for Linux (WSL) for Windows users), as well as running a \lstinline[language=Go]!go install! command, cloning the repository, and running \lstinline[language=Go]!go generate \&\& go build! \\
		Documentation & Contains a Godoc page, wiki page, an examples subdirectory, and a detailed README.md providing what PEG features are supported, command line interface usage, installation, and examples & Contains a detailed README.md and a PEG file syntax file containing what PEG features are supported, installation steps, and example projects links \\
		Generated Code & Generated Go code can be used to configure settings on the parser, view statistics or error information, debug, enable or disable memoization, and parse readers, files, and strings. Additionally, the generated code is well documented, with several comments surrounding usage of public functions & Generated Go code can be used to run the PEG script on a provided string, debug and output the Abstract Syntax Tree (AST), output errors in parsing, and enable or disable memoization \\
		Examples & Provides three direct examples: detecting indentation in files, parsing JSON files, and a calculator example. Examples were kept short (in regards to PEG scripts) with the calculator and indentation examples showcasing usage via a main function & Provides two project links for examples on PEG usage: a natural date/time parser and a scientific names parser. Both projects have large PEG scripts, with most lines being simple to read. The natural date/time project was intuitive to use and understand program flow \\
	\end{tabular}}
\end{table}

While examples are relatively similar, the simplicity of installation for Pigeon as opposed to \acrshort{peg}, combined with the depths of documentation both in generated code and outwith it, makes Pigeon a more ideal option. Additionally, its greater number of features gives LinTeX more to make use of. Therefore, Pigeon was the chosen library for the project.

\subsection{Choice of frameworks and libraries}
\subsubsection{Backend frameworks and libraries}
There were three options for web frameworks and packages, namely Gin \parencite{go-gin}, Gorilla \parencite{gorilla-mux}, and net/http \parencite{net-http}. For \lintex , net/http was chosen for the key reason that it had the least barriers to entry, as there was a provided guide on the Go development website that was utilised in the project setup. However, in hindsight more consideration towards what web framework to use would have had a great impact on the project as a whole, as net/http is a lower level web framework than Gin or Gorilla, consequently requiring more management on \lintex's end. Some examples of difficulties that occurred is a lack of POST requests being easily supported and necessitating server files are explicitly stated as opposed to implicitly defined within a given context.

The html/template \parencite{html-template} package was included in the project. This package provided the ability to load \acrfull{html} files through Go, while allowing information to be provided to them. The alternative choice was to provide a raw print output to the client, which made the package a far more clear winner.

For the project to be able to compile \latexspaced documents to convert them into \acrshortpl{pdf}, it was necessary to either use a pre-built package or to compile \latexspaced documents locally. Pre-built packages that were found include go-latex \parencite{go-latex} and latex-fast-compile \parencite{latex-fast-compile}. It was necessary to rule out go-latex as it emulated a \latexspaced compiler, rather than directly linking to one. Additionally, it lacked documentation and necessary features making it unsuitable. Latex-fast-compile was also unsuitable due to using \pdflatexspaced rather than \lualatexspaced which was desired.

Consequently, it became necessary to compile \latexspaced documents directly. This was done through the use of \acrfull{cmd} which allowed the `lualatex' command to be ran, consequently compiling the requested file. It did, however, prerequisite that the hosting server had a \latexspaced compiler in their backend, for the relevant \lualatexspaced engine. The os/exec \parencite{os-exec} package allowed for a \acrshort{cmd} shell to be created within Go code, consequently allowing the `lualatex' command to be ran. Hence, compiling the \latexspaced document that the user provides.

\subsubsection{Frontend frameworks and libraries}
As plain \acrfull{css} is difficult to work with and maintain, it was necessary to decide on a \acrshort{css} framework for the project. My prior experience in Bootstrap \parencite{bootstrap} made it a close choice, but the opportunity to learn to utilise Tailwind CSS \parencite{tailwind} was more enticing and eventually won the decision.

To implement Tailwind CSS, it was necessary to install its \acrfull{cli} through \acrfull{npm}. This permitted usage of Tailwind CSS without necessitating \acrfull{js} within the codebase, and avoided workarounds such as a \acrfull{cdn}. Additionally, due to difficulties in using Tailwind CSS, it was necessary to utilise a component library built on top of it, for which Flowbite \parencite{flowbite} was chosen as it could be used as almost plug-and-play.

To ensure appropriate responsiveness of \lintex, it was necessary to be able to introduce updates in real time. For this, HTMX \parencite{htmx} was utilised, as it is a tool that permits \acrshort{html} tags other than the `a' and `form' tags to send \acrfull{http} requests to the server. For example, it is possible to use HTMX to create a \acrshort{http} request when a given key is pressed through its event management.

\subsection{Guidelines for accessibility}
The following key points have been used in the design of \lintexspaced to ensure its accessibility, following WCAG 2.2 \parencite{wcag} standards.
\begin{itemize}
\item The website should have all non-text content contain an \gls{alt_text} (WCAG 2.2: Guideline 1.1).
\item The website may be operable through the keyboard alone (WCAG 2.2: Guideline 2.1).
\item LinTeX should provide simple error messages.
\item The website should allow for content to be distinguishable including supporting resizing content and minimal colour usage (WCAG 2.2: Guideline 1.4).
\item The website should be presentable in a simpler format if required (WCAG 2.2: Guideline 1.3).
\item The website should include navigation features (WCAG 2.2: Guideline 2.4).
\item The website should be predictable (WCAG 2.2: Guideline 3.2).
\item The produced document by \lintexspaced should contain appropriate tagging.
\item The produced \gls{accessible_doc} by \lintexspaced should contain appropriate \gls{alt_text}, actual text, or artifact.
\end{itemize}

\subsection{Choosing a software development methodology}
The available options under consideration for the project are the waterfall methodology, agile development methodology, \acrfull{tdd}, and \acrfull{bdd}.

The waterfall methodology requires significant research early in the project to have success. However, as I lack experience in \acrshort{peg} script usage and many related topics for this project, it was likely changes would be needed dynamically.

\Acrshort{tdd} was the next consideration, with a personal strong bias towards it due to its usefulness in being self-documenting and saving the hassle of writing the unit tests after creating the software, which can be a more draining process. However, as Pigeon, the PEG library chosen in \fullref{backg:peg_library_choice}, generates Go files without a way to test individual components, it was not possible to create unit tests without it using complete \latexspaced code for the feature that should be tested. Additionally, creating unit tests can be time-consuming, which may be unsuitable with the short timeframe of the project.

\Acrshort{bdd} was removed from consideration for similar reasons as \acrshort{tdd}, with the addition of it being more suitable for larger contexts such as games development, where it is a lot less suitable to test every feature individual.

The final methodology, agile development methodology, allowed for code to be created quickly through its sprints. Additionally, as there was no prerequisite to create unit tests before creating the features, it saved a lot of time in development for the project. Consequently, agile development methodology was chosen for this project.

\subsection{Screen reader choice}
Through a small number of critera, as showcased in \tableref{table:screen_reader_choices}, it was possible to reduce the choices for what screen reader to use in our context.

\begin{table}[h]
	\caption{Screen reader pricing status and \acrshortpl{os}}\label{table:screen_reader_choices}
    \centering
    \resizebox{\textwidth}{!}{\begin{tabular}{ccc}
    Screen reader & Paid service & \acrlong{os} support \\\hline
    JAWS & Yes & Microsoft Windows 10, 11, and server editions \\
    Apple VoiceOver & No & Apple devices including iPhones, iPads, Macs, and AppleWatches \\
    Windows Narrator & No & Windows 10 and 11 \\
    TalkBack & No & Android 6.0 Marshmallow onwards \\
    ChromeVox & No & Primarily ChromeOS, with older versions supporting devices using Chrome browser \\
    Supernova & Yes & Microsoft Windows 10, 11, server editions, and virtual environments \\
    NVDA & No & Microsoft Windows 8.1 to 11 \\
    \end{tabular}}
\end{table}

As a student, I do not have the leeway to purchase subscriptions for several screen readers, making those that require payments ruled out, even if it would just be in the form of a free trial. Additionally, there is an ethical dilemma in requiring payments for a necessary accessibility tool, which pushed me away from choosing them. Another criteria was that the Windows \acrfull{os} should be supported, with emphasis placed on Windows 10 and Windows 11. These were chosen as they are the most common PC \acrshort{os} used, consequently targetting a larger base of users who may require screen readers.

Considering these critera, NVDA and Microsoft Narrator were the final choices. When testing both on a placeholder document, emulating the experience of a screen reader user working with \latex , I had found that there was more difficulty with Windows Narrator in terms of having it read the correct extract on the document, often requiring using a mouse to get to the correct location. Inversely, NVDA read the correct extract in full detail. Hence, by process of elimination to a degree, NVDA was the chosen screen reader for \lintex.

Adobe Acrobat Reader was used to work alongside NVDA to make the documents easier to navigate in more complex contexts, such as where the document may not be read as simply as top-to-bottom, left-to-right.

\subsection{\nmu User \nmu study design choices}\label{backg:user_study_design}
\subsection{Designing a \acrshort{gui}}


LinTeX will take strong inspiration from Integrated Development Environments (IDEs) with regards to its User Interface/User Experience (UI/UX) design. This is because LaTeX in structure, is similar to a programming language, even being Turing complete \cite{latexturingcomplete}. Hence, an environment suitable to programming is suitable to LaTeX, with some special considerations put aside for its quirks. For example, depending on the user's screen size and preferences, they may wish to have the output of the LaTeX documents taking up half of the screen, a separate browser window or monitor, or out of the way until compilation.

- Tailwind choice?
- Flowbite choice?
- Potential HTMX usage if time allowed for it?